\documentclass[11pt]{article}

\usepackage[margin=1in]{geometry}
\usepackage[T1]{fontenc}
\usepackage[utf8]{inputenc}
\usepackage{lmodern}
\usepackage{amsmath}
\usepackage{booktabs}
\usepackage{tabularx}
\usepackage{enumitem}
\usepackage{xcolor}
\usepackage[hidelinks]{hyperref}

\setlist{nosep, leftmargin=*}
\newcommand{\todo}[1]{\textcolor{red}{[#1]}}

\title{ECSE 552 Project Proposal\\[0.5em]
\large Teaching a Robot Policy a New Object from Action-Free Videos:\\
Object-Invariant Inverse Dynamics for Pseudo-Labeling}
\author{Brik \and Zheng \and Myriam \and Gonzalo}
\date{Draft v0 --- initial outline by Gonzalo, October 5, 2026}

\begin{document}

\maketitle

% ---------------------------------------------------------------------------
\section{Title and Summary}

\begin{itemize}
  \item \textbf{Working title:} Teaching a Robot Policy a New Object from
    Action-Free Videos: Object-Invariant Inverse Dynamics for Pseudo-Labeling.
  \item \textbf{One-line summary:} use action-labeled demos of one object plus
    action-free videos of a second object to train a policy for the second
    object, entirely in simulation.
  \item \textbf{Central question:} where should domain invariance be applied ---
    in the model that labels the videos, in the policy, or both?
\end{itemize}

% ---------------------------------------------------------------------------
\section{Background}

\subsection{Application (A1)}
\begin{itemize}
  \item Learning robot manipulation skills (picking up and handing over
    objects) from demonstrations.
  \item Imitation learning / behavior cloning: a neural network maps camera
    images and joint positions to motor commands.
  \item Current state of the art for this setting: ACT (Action Chunking with
    Transformers)~\cite{zhao2023act}.
\end{itemize}

\subsection{Problem (A2)}
\begin{itemize}
  \item Behavior cloning needs demonstrations where every frame is paired with
    the motor command used.
  \item These labeled demos are slow and expensive to collect (a human must
    teleoperate the robot).
  \item Videos of a robot doing a task are much cheaper, but contain no motor
    commands.
  \item A policy trained on one object fails on a new object that needs a
    different grasp.
  \item \textbf{Question:} can action-free videos of a new object, plus labeled
    demos of a known object, train a policy for the new object?
\end{itemize}

\subsection{Why a Solution Is Needed (A2)}
\begin{itemize}
  \item Collecting labeled demos for every new object does not scale.
  \item Prior domain adaptation work~\cite{cheng2025gda} assumes the task is
    the same across domains and only appearance changes (e.g.\ the color of the
    object).
  \item Here the correct actions change with the object, so standard
    adversarial alignment may remove information the policy needs.
  \item Pseudo-labeling videos with an inverse dynamics model works when the
    labeler generalizes; whether it survives an object change is untested.
\end{itemize}

\subsection{Constraints (A3)}
\begin{itemize}
  \item \textbf{No target action labels:} only image frames of the new object
    may be used for training.
  \item \textbf{Closed-loop task success:} the goal is a high success rate when
    the policy controls the simulated robot.
  \item \textbf{Real-time inference:} the policy must output actions within the
    50\,Hz control budget ($<20$\,ms per step).
  \item \textbf{Data efficiency:} at most $\sim$50 labeled source
    demonstrations.
  \item \textbf{Selective invariance:} the labeling model should ignore object
    identity; the policy must stay sensitive to it.
  \item \textbf{Robustness to compounding error:} small action errors must not
    drift the robot off the training distribution.
\end{itemize}

% ---------------------------------------------------------------------------
\section{Aims and Goals}

\subsection{Aims}
\begin{enumerate}[label=\textbf{Aim \arabic*:}, leftmargin=4.2em]
  \item Build an inverse dynamics model (IDM) that infers motor commands from
    consecutive frames and stays accurate when the object changes.
  \item Use it to pseudo-label action-free videos of a new object, then train
    an ACT policy on labeled + pseudo-labeled data.
  \item Test where invariance belongs: object-invariant IDM vs.\
    object-invariant policy vs.\ both.
\end{enumerate}

\subsection{Scope (A4)}
\begin{itemize}
  \item \textbf{In scope:} one source task (ALOHA Transfer Cube), one target
    object (elongated bar at random rotation), simulation only.
  \item \textbf{In scope:} evaluation of the full pipeline in closed-loop
    simulated rollouts.
  \item \textbf{Out of scope:} real robots, human videos, other robot arms,
    multiple tasks, language instructions.
\end{itemize}

\subsection{Proposed Solution (A5)}
\begin{itemize}
  \item \textbf{Key insight:} in ALOHA an action is the next target joint
    position, so the action is visible in the next frame.
  \item An IDM reading frames $t$ and $t+1$ is mostly estimating arm pose,
    which does not depend on the object.
  \item Therefore, make the IDM object-invariant with an adversarial domain
    classifier and a gradient reversal layer.
  \item Do \emph{not} make the policy invariant, since it must grasp each
    object differently.
  \item \textbf{Pipeline:} train IDM on cube data $\rightarrow$ label bar
    videos $\rightarrow$ train ACT on both.
\end{itemize}

% ---------------------------------------------------------------------------
\section{Methodology}

\subsection{Data Generation (A6)}
\begin{itemize}
  \item \textbf{Simulator:} ALOHA bimanual environment in
    MuJoCo~\cite{todorov2012mujoco}, via LeRobot~\cite{cadene2024lerobot}
    \texttt{gym-aloha} and the ACT repository.
  \item \textbf{Source data:} $\sim$50 scripted Transfer Cube demos with images
    and actions.
  \item \textbf{Target data:} 10--100 scripted demos with a new elongated bar;
    the scripted policy is extended to rotate the gripper.
  \item Target actions are hidden from all models; they are kept only for
    evaluation and for the oracle.
  \item \textbf{Held-out test set:} new random object positions, never used for
    training.
\end{itemize}

\subsection{Models and Implementation (A5, A6)}
\begin{itemize}
  \item \textbf{Policy:} ACT --- ResNet-18 image encoder, transformer
    encoder--decoder, CVAE objective, action chunks.
  \item \textbf{IDM:} ResNet-18 encoder on frames $(t, t+1)$ $\rightarrow$ MLP
    $\rightarrow$ 14-dim action; trained with L1 loss on cube data.
  \item \textbf{Object-invariant IDM:} add a domain classifier on IDM features
    (cube vs.\ bar frames) with a gradient reversal layer.
  \item \textbf{Training loss (IDM):}
    \begin{equation*}
      \mathcal{L}_{\text{IDM}} = \mathcal{L}_{1}^{\text{action}}
      + \lambda \, \mathcal{L}_{\text{adv}}^{\text{domain}},
    \end{equation*}
    with $\lambda$ tuned on a validation split.
  \item \textbf{Final policy:} ACT trained on cube demos + pseudo-labeled bar
    demos.
  \item \textbf{Framework:} PyTorch + LeRobot; hyperparameters tuned on a
    validation set of held-out episodes.
\end{itemize}

\subsection{Evaluation (A7)}
\begin{itemize}
  \item \textbf{Primary metric:} closed-loop success rate on the bar,
    50 rollouts $\times$ 3 seeds per method.
  \item Subtask success: grasp, lift, handover.
  \item Pseudo-label quality: IDM action error vs.\ hidden true actions.
  \item Inference latency per step (checks the real-time constraint).
  \item Scaling study: success vs.\ number of target videos
    (10, 25, 50, 100).
  \item Ablation: invariance in IDM only / policy only / both.
  \item t-SNE plots of cube vs.\ bar features in the IDM and in the policy.
\end{itemize}

\subsection{Alternative Solutions, All Run on Our Data (A8)}
\begin{itemize}
  \item \textbf{Source-only ACT:} trained on cube data only (lower bound).
  \item \textbf{ACT + DANN:} adversarial invariance on the policy encoder,
    using unlabeled bar frames.
  \item \textbf{IDM pseudo-labeling without adaptation:} VPT / BCO-style.
  \item \textbf{Diffusion Policy:} alternative state-of-the-art policy
    backbone, trained with the best pipeline.
  \item \textbf{Oracle ACT:} trained on true bar actions (upper bound).
\end{itemize}

\subsection{Why Ours Could Be Superior (A8)}
\begin{itemize}
  \item DANN on the policy forces cube and bar features together, discarding
    the grasp information the bar needs.
  \item An unadapted IDM may mislabel frames when it sees an unfamiliar object.
  \item Our method applies invariance only where the input--output mapping is
    shared across objects.
\end{itemize}

% ---------------------------------------------------------------------------
\section{Feasibility and Resources}
\begin{itemize}
  \item \textbf{No robot needed:} all data generation and evaluation run in
    the MuJoCo simulator, which is open-source.
  \item \textbf{Dataset:} self-generated with scripted policies; ground-truth
    actions are available for every frame, so pseudo-labels can be scored
    exactly.
  \item \textbf{Code:} open-source ACT repository (scripted policies, episode
    recording) and LeRobot (ACT, Diffusion Policy, \texttt{gym-aloha}).
  \item \textbf{Compute:} ACT ($\sim$80M parameters) trains in $\sim$5\,h on
    one RTX 2080 Ti in the original paper~\cite{zhao2023act}.
  \item Estimated $\sim$15--20 training runs $\rightarrow$ reduce cost with
    1--2 cameras instead of 4, lower image resolution, and fewer training
    steps.
  \item \textbf{GPU source:} \todo{TO CONFIRM --- university cluster / Colab
    Pro / personal GPU}.
  \item \textbf{Early check:} confirm the scripted cube policy runs and the new
    bar object loads before finalizing the plan.
\end{itemize}

% ---------------------------------------------------------------------------
\section{Risk Analysis}

\begin{center}
\renewcommand{\arraystretch}{1.3}
\begin{tabularx}{\linewidth}{@{}>{\raggedright\arraybackslash}p{0.38\linewidth}X@{}}
\toprule
\textbf{Risk} & \textbf{Mitigation / fallback} \\
\midrule
Adding the bar or extending the scripted policy is harder than expected
  & Use a visual-only object change (color/texture), or the existing
    Insertion task as target. \\
Pseudo-labels too noisy for the policy to learn from
  & Smooth labels over action chunks; filter low-confidence labels; few-shot
    variant with 5 truly labeled bar demos. \\
Even the oracle policy fails on the bar
  & Simplify the object (shorter bar, smaller rotation range); report
    per-subtask success. \\
Adversarial training unstable
  & Tune $\lambda$, warm-up schedule for the domain loss; try MMD or OT
    alignment as alternatives. \\
Not enough GPU time
  & Smaller models, fewer seeds, offline action error as secondary
    evidence. \\
Results only work in simulation
  & State as a limitation; sim-to-real transfer listed as future work. \\
\bottomrule
\end{tabularx}
\end{center}

% ---------------------------------------------------------------------------
\section{Role of Group Members}

\begin{center}
\renewcommand{\arraystretch}{1.3}
\begin{tabularx}{\linewidth}{@{}lX@{}}
\toprule
\textbf{Member} & \textbf{Main responsibilities} \\
\midrule
Brik    & \todo{TBD} \\
Zheng   & \todo{TBD} \\
Myriam  & \todo{TBD} \\
Gonzalo & \todo{TBD} \\
\bottomrule
\end{tabularx}
\end{center}

\noindent All members: literature review, proposal, presentation, and report
sections for their component.

% ---------------------------------------------------------------------------
\bibliographystyle{unsrt}
\bibliography{../../references}

\end{document}
